\begin{abcliste}
\abc The charge stays on the plates, so use $E = \frac{Q}{\varepsilon_0 A}$: the field depends on the charge per area, not on the distance.
\begin{align}
 \frac{E_2}{E_1} &= \frac{2Q/(\varepsilon_0 A/2)}{Q/(\varepsilon_0 A)} = \frac{2}{1/2} = \result{4}
\end{align}

\abc $F = |q|\,E$: the field is four times as strong and the alpha particle carries twice the charge of the electron.
\begin{align}
 \frac{F_2}{F_1} &= \frac{2e\cdot E_2}{e\cdot E_1} = 2\cdot 4 = \result{8}
\end{align}
The masses do not matter for the force.

\abc In both capacitors the field points down, from the positive to the negative plate. The force on the positive alpha particle points down, the force on the negative electron up: \result{\text{opposite ways}}.

\abc Now the sources keep the voltages, so use $E = \frac{U}{d}$: the field depends on the voltage and the distance, not on the area.
\begin{align}
 \frac{E_2}{E_1} &= \frac{2U/(2d)}{U/d} = \result{1}, &
 \frac{F_2}{F_1} &= \frac{2e\cdot E_2}{e\cdot E_1} = \result{2}
\end{align}
\end{abcliste}
